\documentclass[12pt]{article}

\usepackage[margin=1in]{geometry}
\usepackage{amsmath}
\usepackage{amssymb}
\usepackage{graphicx}
\usepackage{setspace}
\usepackage{microtype}
\usepackage[T1]{fontenc}
\usepackage{newtxtext}
\usepackage{newtxmath}
\usepackage[hidelinks]{hyperref}

\doublespacing
\setlength{\parindent}{0.5in}
\setlength{\parskip}{0pt}
\setlength{\emergencystretch}{3em}
\setcounter{secnumdepth}{3}
\setcounter{tocdepth}{2}

\title{Physically Grounded Camera Based Estimation of\\
Skin Reflectance, Illumination, Physiology, and Affect\\
Across Skin Tones}
\author{Emily Rooney}
\date{July 2026}

\begin{document}

\begin{titlepage}
\centering
\vspace*{1in}

{\Large\bfseries
Physically Grounded Camera Based Estimation of\\
Skin Reflectance, Illumination, Physiology, and Affect\\
Across Skin Tones\par}

\vspace{1in}
{\large Emily Rooney\par}
\vspace{0.5in}

Program of Color Science\\
Rochester Institute of Technology

\vfill

Research collaborators\\
Mike Murdoch\\
Chris Thorstenson\\
Kaida Xiao\\
Ferat Sahin

\vfill
July 2026
\end{titlepage}

\pagenumbering{roman}
\tableofcontents
\clearpage
\pagenumbering{arabic}

\section{Introduction}

Human skin contains measurable information about surface appearance, tissue composition, cardiovascular activity, and affect related behavior. Cameras provide a practical and noncontact means of observing these signals. Camera based measurement is therefore relevant to color science, physiological monitoring, affective computing, and human computer interaction. A camera does not measure these properties directly. Its recorded values result from interactions among the spectral power distribution of the illumination, the spectral reflectance of the skin, tissue dependent light absorption and scattering, camera spectral sensitivities, exposure settings, and subsequent image processing operations.

Several signals can coexist within the same facial pixels. Stable spatial differences may describe skin reflectance. Temporal color changes may encode blood volume variation. Illumination changes can alter both types of information. Facial movement and expression may provide evidence related to affect. These sources of variation can reinforce or confound one another. A model trained directly on rendered RGB imagery may learn useful statistical relationships without separating the physical factors responsible for the recorded signal. This can reduce its reliability when it is applied to new cameras, lighting conditions, environments, or populations.

This dissertation approaches facial imaging as a physically grounded measurement problem. Its four projects examine how color representation affects affect estimation, how skin color can be recovered from flash and no flash images, how scene illumination can be estimated from facial appearance, and whether camera derived physiological information improves valence and arousal estimation. Together, these projects investigate when colorimetric modeling can improve the accuracy, interpretability, generalizability, and equity of camera based measurements derived from human skin.

\section{Background}

\subsection{Camera Based Measurement and Image Formation}

A digital camera does not directly measure skin reflectance, illumination, physiological activity, or affect. It records channel responses produced by the interaction among incident illumination, surface and subsurface optical properties, camera spectral sensitivities, exposure, and image processing. Interpreting facial imagery as measurement data therefore requires a model of how these factors contribute to the recorded signal.

For a linearized camera response, the value recorded by channel \(c\) at image location \(x\) can be approximated as

\begin{equation}
I_c(x)
=
g\int_{\lambda}E(\lambda,x)R(\lambda,x)S_c(\lambda)\,d\lambda
+
\epsilon_c(x),
\end{equation}

where \(E(\lambda,x)\) is the spectral power distribution of the illumination, \(R(\lambda,x)\) is the spectral reflectance of the observed surface, \(S_c(\lambda)\) is the spectral sensitivity of camera channel \(c\), \(g\) represents exposure dependent scaling, and \(\epsilon_c(x)\) represents noise and modeling error.

An RGB camera reduces the continuous spectral signal to three channel values. Different combinations of illumination and reflectance can therefore produce similar RGB responses. Recovery of either quantity from a single image is generally underdetermined unless acquisition includes calibration data, additional observations, or justified physical and statistical constraints.

The linear model is most applicable to characterized or linearized image data. Rendered JPEG and video data may include demosaicing, white balance, denoising, color correction, tone mapping, sharpening, gamma encoding, and compression. These operations alter the relationship between scene radiance and pixel values. RAW data can preserve a closer relationship to sensor measurements, although computational RAW formats such as Apple ProRAW may still include demosaicing, multiple exposure fusion, gain maps, and other processing.

In facial video, illumination, reflectance, motion, physiology, and camera settings may all vary over time. Stable spatial information can support estimation of skin appearance. Temporal modulation caused by blood volume changes can support physiological estimation. Facial configuration and movement can support affect prediction. These quantities share the same recorded pixels but require different assumptions, processing methods, and reference measurements.

\subsection{Camera Colorimetry and Color Representation}

Camera RGB values depend on both the scene and the imaging system. The same skin may produce different RGB values under different illuminants, cameras, exposure settings, or processing pipelines. Camera RGB should therefore not be interpreted as a direct measurement of surface color.

Camera spectral sensitivities do not generally match the CIE color matching functions used to describe human observer colorimetry. Under the Luther Ives condition, camera spectral sensitivities must be expressible as a linear transformation of the CIE color matching functions for exact colorimetric recovery across arbitrary spectra. Consumer cameras generally do not satisfy this condition.

A transformation from linear camera RGB to CIE XYZ can nevertheless be fitted for a restricted set of reflectances and acquisition conditions:

\begin{equation}
\begin{bmatrix}X\\Y\\Z\end{bmatrix}
=
\mathbf{M}
\begin{bmatrix}R\\G\\B\end{bmatrix},
\end{equation}

where \(\mathbf{M}\) is a fitted \(3 \times 3\) color correction matrix. These transformations can be estimated from images of samples with known reference values, such as a ColorChecker. Calibration focused on skin colored samples may provide greater accuracy within the restricted skin color region than a general chart fit \cite{harville2005,choi2017}.

The accuracy of a fitted transformation depends on the camera, illuminant, exposure, processing pipeline, and reflectances represented during calibration. Offline camera characterization can remove the need for a target during later use, but it assumes that the imaging pipeline and deployment conditions remain sufficiently stable.

Recovery of CIE tristimulus values should be distinguished from recovery of spectral reflectance. A color correction matrix estimates three dimensional XYZ values under the conditions represented by its calibration data. It does not reconstruct the full wavelength dependent reflectance of the surface.

Camera RGB is device dependent because its values depend on sensor sensitivities and processing. CIE XYZ is device independent when derived from valid colorimetric measurements. CIELAB is calculated from XYZ relative to a specified reference white and represents lightness through \(L^*\) and approximately opponent chromatic dimensions through \(a^*\) and \(b^*\). Colorimetric error can be described with coordinate differences or formulas such as CIEDE2000 \cite{sharma2005}.

YCbCr separates a luma related component from two chroma related components and is widely used in digital video. Unlike XYZ and CIELAB, it is not intended as a device independent colorimetric representation. Color space conversion reorganizes image information but does not create new information. It may nevertheless affect a learning model because it changes channel distributions and relationships.

\subsection{Human Skin Optics, Color Measurement, and Equity}

Human skin is a layered and heterogeneous optical material. Its appearance results from surface reflection, absorption, scattering, and light transport within tissue. Melanin within the epidermis is a major contributor to visible pigmentation. Hemoglobin within the dermis also contributes wavelength dependent absorption, while changes in blood volume produce small temporal variations in reflected light.

Skin appearance cannot be described by melanin alone. Hemoglobin, tissue structure, hydration, sebum, cosmetics, surface texture, and subsurface scattering also affect the measured signal. Skin is not a purely diffuse surface. Specular reflection may preserve characteristics of the illumination, while subsurface reflection is more strongly modified by tissue absorption and scattering.

\subsubsection{Measurement Geometry and Region Selection}

Measured skin color depends on the relative positions of the camera, participant, and light source. Facial curvature changes the angle of incidence and observation across facial locations. Directional illumination may produce highlights, shadows, and spatially varying shading.

Contact instruments and cameras use different measurement geometries. A contact colorimeter measures a small region under controlled illumination, while a camera records a larger area that may include gloss, pores, hair, shadows, and spatial variation. Agreement must therefore be evaluated using clearly defined and repeatable regions.

Cheek regions are often suitable for facial color measurement. Regions containing facial hair, eyes, lips, strong highlights, deep shadows, or highly curved boundaries should be excluded. Facial landmarks or mesh models can support repeatable region selection.

\subsubsection{Defining and Measuring Skin Color}

The term skin tone may refer to perceived complexion, measured lightness, categorical labels, pigmentation, or combinations of color attributes. These concepts should not be treated as interchangeable.

The Fitzpatrick Skin Type scale classifies skin according to reported responses to ultraviolet exposure. It was not designed as an objective color measurement scale and is an imperfect proxy for visible skin color \cite{gupta2019}. Visual scales such as the Monk Skin Tone scale provide broader appearance categories but remain ordinal.

Instrumentally measured CIE coordinates provide a stronger basis for continuous analysis. Skin can be described using \(L^*\), \(a^*\), and \(b^*\), or related measures of lightness, chroma, and hue. Contact based cheek measurements can provide reference values for evaluating camera estimates.

\subsubsection{Measurement Equity}

Camera performance may vary across skin reflectances because the number and spectral distribution of detected photons differ across participants and wavelengths. Lower reflected intensity can reduce signal to noise ratio. Automatic exposure and tone processing may compensate for some variation but may also introduce nonlinear processing.

Performance differences should be evaluated using continuously measured skin color whenever possible. Race, ancestry, Fitzpatrick type, visual skin tone categories, and measured skin lightness describe different constructs and should not be used as substitutes for one another.

Colorimetric analyses should report overall color difference together with errors in \(L^*\), \(a^*\), and \(b^*\). Physiological analyses should report accuracy together with signal coverage and failure rate.

\subsection{Project One: Color Representation for Affect Estimation}

Project One examines whether the color representation supplied to a deep learning model affects continuous valence and arousal estimation. This question concerns the encoding of visual input rather than the representation of emotion itself.

\subsubsection{Continuous Valence and Arousal Estimation}

Emotion involves subjective experience, physiological activity, expression, appraisal, and behavior. Facial images provide observable evidence about expression and behavior but do not provide direct access to internal emotional experience.

Continuous affect models commonly represent affect through valence and arousal. Valence describes the positive or negative character of an affective state, while arousal describes its level of activation \cite{russell1980}. Observer annotations describe how an expression appears to another person, while participant self reports describe the experience reported by the participant.

\subsubsection{EmoNet and RGB Pretraining}

AffectNet provides facial images annotated for categorical expression and continuous valence and arousal \cite{mollahosseini2017}. AFEW VA extends continuous affect analysis to naturalistic video \cite{kossaifi2017}. EmoNet uses a Face Alignment Network backbone and learned attention for continuous affect estimation.

Because its feature extractor was developed with RGB imagery, alternative representations supply channel statistics outside its original training distribution. Performance differences may therefore reflect compatibility with RGB pretrained features rather than an intrinsic advantage of one color space.

\subsubsection{Color Representation in Deep Learning}

RGB is commonly used because public datasets and pretrained networks are distributed in that form. Previous work shows that convolutional neural networks are sensitive to color encoding and that color space selection can affect computer vision tasks \cite{singh2020assessing,gowda2018colornet}. Color normalization has also affected facial expression recognition \cite{nunnari2025color}.

Changing color space alters how luminance and chromatic information are distributed among channels. A controlled comparison must hold architecture, initialization, data partitions, training protocol, and optimization constant while changing only the input representation.

\subsubsection{Domain Shift and Project Gap}

Differences in camera, compression, illumination, pose, participants, context, and annotation can produce performance degradation on new data. SEWA contains naturalistic interactions with continuous affect annotations \cite{kossaifi2019}. MABL ME instead uses trial level participant self reports.

Project One compares RGB, XYZ, CIELAB, and YCbCr using a frozen RGB pretrained EmoNet backbone and evaluates whether the observed relationship persists on unseen datasets. Its completed methodology and findings are presented in the Project One section of this proposal.

\subsection{Project Two: Camera Based Skin Reflectance Estimation}

Camera based skin reflectance estimation seeks to recover stable information about skin appearance while reducing the effects of illumination, geometry, camera response, and processing. A visually plausible skin image is not necessarily a colorimetrically accurate measurement.

\subsubsection{Target Based and Offline Characterization}

Physical color targets support traceable camera characterization but are impractical in every consumer capture. Offline characterization fits a transformation from controlled target captures for later target free use. This assumes that camera hardware, processing, and deployment conditions remain sufficiently stable.

\subsubsection{Flash and No Flash Capture}

Flash and no flash imaging provides two observations under related illumination. The ambient image records existing illumination, while the flash image records ambient and flash contributions. Under a linear response, their difference provides information about the flash contribution \cite{dicarlo2001,petschnigg2004,lu2006,hui2016,hui2018,maralan2023}.

Reliable subtraction requires matched camera settings and accurate alignment. Movement, blinking, expression changes, saturation, specular highlights, and spatially nonuniform flash illumination can introduce error. A flash pair does not automatically provide absolute reflectance.

\subsubsection{Intrinsic Decomposition and Inverse Rendering}

Intrinsic decomposition separates reflectance from illumination dependent shading. Inverse rendering may additionally estimate geometry, lighting, camera properties, and material parameters \cite{garces2022}. Human skin challenges simple Lambertian assumptions because it exhibits specular reflection, subsurface transport, spatial variation, and wavelength dependent absorption.

CGIntrinsics demonstrated the value of physically based synthetic supervision \cite{li2018cgintrinsics}. Mobile reflectometry and two image material capture provide precedents for commodity camera and flash acquisition \cite{riviere2016,aittala2015}. Portrait relighting models more directly account for facial geometry, specular reflection, and shadow \cite{nestmeyer2020,wang2020portrait}. These methods generally optimize appearance rather than absolute colorimetric agreement.

\subsubsection{Reflectance Research Gap}

There has been limited evaluation of target free flash and no flash skin color estimation against contact based measurements across a broad skin color range. Project Two addresses this gap through offline characterization, paired capture, defined facial regions, and contact based CIELAB reference measurements.

\subsection{Project Three: Facial Illumination Estimation}

Illumination estimation seeks to determine properties of the light reaching a scene. Skin pigmentation and illumination color can produce related changes in camera RGB, creating a confound for facial imagery.

\subsubsection{Computational Color Constancy}

Classical color constancy methods rely on assumptions about scene reflectance. Gray world assumes that average reflectance is neutral, white patch assumes that the brightest region is white, and gray edge uses assumptions about spatial reflectance changes \cite{gijsenij2011}. These assumptions may fail in portraits that contain no neutral surface.

Learning based methods estimate illumination from annotated images but may depend on the cameras, illuminants, and scenes represented during training \cite{xie2023}. A facial model may also learn associations between skin appearance and illumination. Supplying measured skin color as context may help, but it can introduce circular reasoning if estimated from the same uncorrected image.

\subsubsection{Illuminant Representation and Evaluation}

Camera white balance commonly represents illumination through RGB gains or camera chromaticity. With valid XYZ values, illumination may be represented using CIE chromaticity. Correlated color temperature describes the closest point on the Planckian locus, while \(D_{uv}\) describes signed distance from that locus. Neither quantity recovers the complete spectral power distribution.

Predictions should be compared with measurements from a calibrated colorimeter or spectroradiometer. Evaluation may include angular error, chromaticity error, correlated color temperature error, and \(D_{uv}\) error. The practical effect can also be tested by correcting facial images and comparing estimated skin color with contact measurements.

\subsubsection{Illumination Research Gap}

General color constancy methods are not designed specifically for colorimetric analysis of human skin. Project Three estimates physically interpretable illumination quantities from facial imagery and tests whether measured skin color reduces confusion between pigmentation and illumination.

\subsection{Project Four: Physiological Enhancement of Affect Estimation}

% TODO: Replace this placeholder with the condensed physiological and
% multimodal background in the next editing pass.

\subsection{Background Synthesis and Dissertation Gap}

The four projects share a common camera measurement framework but address distinct inference problems. Study One examines visual input representation. Study Two estimates stable skin color. Study Three estimates illumination. Study Four examines temporal physiological variation and its value for affect prediction.

Physically meaningful recovery requires controlled acquisition, camera characterization, additional observations, or justified priors. Performance on one dataset does not establish recovery of the intended physical quantity rather than a dataset specific correlation. The proposed projects therefore rely on independent reference measurements and explicit evaluation across skin colors, participants, sessions, and domains.

\section{Dissertation Research Questions}

\begin{enumerate}
\item How does input color representation affect continuous valence and arousal estimation under a controlled architecture and training protocol, and does the observed ranking transfer to unseen datasets?
\item Can skin CIELAB values be recovered from matched flash and no flash consumer camera images without a physical target during inference, and how does performance vary across skin colors?
\item Can scene illumination be estimated from a facial image using physically interpretable outputs, and does explicit skin color information reduce confusion between pigmentation and illumination?
\item Does validated camera derived physiological information improve valence and arousal estimation beyond facial video alone?
\end{enumerate}

\section{Proposed Dissertation Projects}

\subsection{Project One: Color Representation for Affect Estimation}
% TODO: Condense the completed conference paper into status, methods,
% results, limitations, and dissertation contribution.

\subsection{Project Two: Camera Based Skin Reflectance Estimation}
% TODO: Add completed work, remaining experiments, validation, and analysis.

\subsection{Project Three: Facial Illumination Estimation}
% TODO: Add dataset collection, model design, validation, and analysis.

\subsection{Project Four: Physiological Enhancement of Affect Estimation}
% TODO: Add physiological validation, multimodal model, and analysis.

\begin{thebibliography}{99}

\bibitem{dicarlo2001}
J. M. DiCarlo, F. Xiao, and B. A. Wandell.
Illuminating illumination.
In \textit{Color and Imaging Conference}, pages 27 to 34, 2001.

\bibitem{petschnigg2004}
G. Petschnigg, R. Szeliski, M. Agrawala, M. Cohen, H. Hoppe, and K. Toyama.
Digital photography with flash and no flash image pairs.
\textit{ACM Transactions on Graphics}, 23, pages 664 to 672, 2004.

\bibitem{lu2006}
C. Lu and M. S. Drew.
Practical scene illuminant estimation via flash and no flash pairs.
In \textit{Color and Imaging Conference}, pages 84 to 89, 2006.

\bibitem{hui2016}
Z. Hui, A. C. Sankaranarayanan, K. Sunkavalli, and S. Hadap.
White balance under mixed illumination using flash photography.
In \textit{IEEE International Conference on Computational Photography}, 2016.

\bibitem{hui2018}
Z. Hui, K. Sunkavalli, S. Hadap, and A. C. Sankaranarayanan.
Illuminant spectra based source separation using flash photography.
In \textit{IEEE Conference on Computer Vision and Pattern Recognition}, pages 6209 to 6218, 2018.

\bibitem{maralan2023}
S. S. Maralan, C. Careaga, and Y. Aksoy.
Computational flash photography through intrinsics.
In \textit{IEEE Conference on Computer Vision and Pattern Recognition}, pages 16654 to 16662, 2023.

\bibitem{harville2005}
M. Harville, H. Baker, N. Bhatti, and S. Susstrunk.
Consistent image based measurement and classification of skin color.
In \textit{IEEE International Conference on Image Processing}, 2005.

\bibitem{choi2017}
H. Choi, K. Choi, and H. J. Suk.
Performance of the 14 skin colored patches in accurately estimating human skin color.
\textit{Electronic Imaging}, 29, pages 62 to 65, 2017.

\bibitem{sharma2005}
G. Sharma, W. Wu, and E. N. Dalal.
The CIEDE2000 color difference formula.
\textit{Color Research and Application}, 30, pages 21 to 30, 2005.

\bibitem{gupta2019}
V. Gupta and V. K. Sharma.
Skin typing with Fitzpatrick grading and other approaches.
\textit{Clinics in Dermatology}, 37, pages 430 to 436, 2019.

\bibitem{gijsenij2011}
A. Gijsenij, T. Gevers, and J. van de Weijer.
Computational color constancy with survey and experiments.
\textit{IEEE Transactions on Image Processing}, 20, pages 2475 to 2489, 2011.

\bibitem{xie2023}
M. Xie, P. Sun, Y. Lang, and M. Fang.
Camera independent color constancy by scene semantics.
\textit{Pattern Recognition Letters}, 171, pages 106 to 115, 2023.

\bibitem{russell1980}
J. A. Russell.
A circumplex model of affect.
\textit{Journal of Personality and Social Psychology}, 39, pages 1161 to 1178, 1980.

\bibitem{mollahosseini2017}
A. Mollahosseini, B. Hasani, and M. H. Mahoor.
AffectNet with a database for facial expression, valence, and arousal computing in natural images.
\textit{IEEE Transactions on Affective Computing}, 10, pages 18 to 31, 2019.

\bibitem{kossaifi2017}
J. Kossaifi and colleagues.
AFEW VA with a database for valence and arousal estimation in natural conditions.
\textit{Image and Vision Computing}, 65, pages 23 to 36, 2017.

\bibitem{kossaifi2019}
J. Kossaifi and colleagues.
SEWA DB with a rich database for audiovisual emotion and sentiment research in natural conditions.
\textit{IEEE Transactions on Pattern Analysis and Machine Intelligence}, 43, pages 1022 to 1040, 2021.

\bibitem{singh2020assessing}
A. Singh, A. Bay, and A. Mirabile.
Assessing the importance of colours for CNNs in object recognition.
\textit{arXiv preprint arXiv:2012.06917}, 2020.

\bibitem{gowda2018colornet}
S. N. Gowda and C. Yuan.
ColorNet: Investigating the importance of color spaces for image classification.
In \textit{Asian Conference on Computer Vision}, pages 581 to 596, 2018.

\bibitem{nunnari2025color}
F. Nunnari, A. Singh, and P. Gebhard.
Color histogram equalization and fine tuning to improve expression recognition of partially occluded faces.
In \textit{Adjunct Proceedings of the ACM International Conference on Intelligent Virtual Agents}, 2025.

\bibitem{garces2022}
E. Garces, C. Rodriguez Pardo, D. Casas, and J. Lopez Moreno.
A Survey on Intrinsic Images: Delving Deep Into Lambert and Beyond.
\textit{International Journal of Computer Vision}, 130, pages 836 to 868, 2022.

\bibitem{li2018cgintrinsics}
Z. Li and N. Snavely.
CGIntrinsics: Better Intrinsic Image Decomposition through Physically Based Rendering.
In \textit{European Conference on Computer Vision}, pages 371 to 387, 2018.

\bibitem{riviere2016}
J. Riviere, P. Peers, and A. Ghosh.
Mobile Surface Reflectometry.
\textit{Computer Graphics Forum}, 35, pages 191 to 202, 2016.

\bibitem{aittala2015}
M. Aittala, T. Weyrich, and J. Lehtinen.
Two Shot SVBRDF Capture for Stationary Materials.
\textit{ACM Transactions on Graphics}, 34, article 110, 2015.

\bibitem{nestmeyer2020}
T. Nestmeyer, J. F. Lalonde, I. Matthews, and A. Lehrmann.
Learning Physics Guided Face Relighting Under Directional Light.
In \textit{IEEE Conference on Computer Vision and Pattern Recognition}, pages 5124 to 5133, 2020.

\bibitem{wang2020portrait}
Z. Wang, X. Yu, M. Lu, Q. Wang, C. Qian, and F. Xu.
Single Image Portrait Relighting via Explicit Multiple Reflectance Channel Modeling.
\textit{ACM Transactions on Graphics}, 39, article 220, 2020.

\end{thebibliography}

\end{document}
